% =============================================================================
%  THE CENTRAL BANK: A SPREAD-CAP BACKSTOP
%
%  INSERT: replaces \subsection{The Central Bank} (\label{sec:cb}) of
%          02-model.tex, i.e. the retired stochastic lending facility
%          (pi^ell, ell_X, cb_t and eqs. P-state-cb ... expectation-compound).
%          Cites eq. hm-yield and eq. spread-def from 02d-sovereign-spread.tex,
%          which comes earlier in the paper (after Endogenous Premiums).
%          A second block at the end replaces eq. (state) and the RCE
%          definition, which still refer to the compound (d, cb) regime.
%
%  Written against global_projection/ on 2026-10-07: point_map.py
%  (_tpi_floor, _tpi_purchase_and_price, _government, _balance_sheets,
%  _incentive_constraint, _tpi_rule), tpi_experiment.py,
%  output_decomposition.decompose_bond_price, calibration/global_projection.py
%  (tpi_cap_bp = 200, tpi_key_D = 0.071, tpi_eps = 1e-4).
%  Notation and macros (\E, \D, \F) follow 02-model.tex.
%
%  NEW SYMBOLS (none collide with 02-model.tex; Z is TFP and m/M index the
%  Gauss-Hermite nodes there, hence the superscripts):
%    \bar c            spread cap, per quarter (200 bp/yr = 0.005)
%    y(q)              HM flow yield
%    \underline q_t    price floor implied by the cap
%    m^{cb}_t          purchases          M^{cb}_t   central-bank book
%    \mathrm{res}_t    reserves held by D banks
%    O^{cb}_t          obligation on reserves issued at t-1
%    \Pi^{cb}_t        central-bank profit (remitted)
%    \kappa_D          capital key, D share
%    \mathcal F_{D,t}  D-bank D-bond first-order gap
%    \tilde q_t        price that clears with no purchase at t
%    \chi_t            Garcia-Zangwill variable (computation paragraph only)
%
%  BIB KEYS REQUIRED:
%    @book{garcia1981pathways, author={Garcia, C. B. and Zangwill, W. I.},
%      title={Pathways to Solutions, Fixed Points, and Equilibria},
%      publisher={Prentice-Hall}, year={1981}}
%    @article{fischer1992special, author={Fischer, A.},
%      title={A special {Newton}-type optimization method},
%      journal={Optimization}, volume={24}, pages={269--284}, year={1992}}
%
%  NUMBERS in "The backstop at work": results/GLOBAL/run.log of 2026-10-06
%    (S_REFINE = 5, 115 points, cap 200 bp/yr, shock p^d -> 2%/qtr). That run
%    PREDATES the uncommitted p'/p correction to the cross-border bond FOCs in
%    point_map._bond_demands; refresh every number after the next full run.
%    "Credit spread" there is the code's 4e4*lambda*mu/varphi.
% =============================================================================


\subsection{The Central Bank: a spread-cap backstop}
\label{sec:cb}

Sovereign risk reaches the real economy in this model through the price of the
periphery bond. A rise in $\pi^{d}_t$ lowers $q^{D}_t$; the fall is a
mark-to-market loss for the banks that hold the bond, which shrinks $n_{D,t}$
through \eqref{eq:ng}, raises $\mu_{D,t}$ through \eqref{eq:mu-closed} and the
credit spread through \eqref{eq:rwc}, and lowers hours through
\eqref{eq:labour-demand}. The policy we study cuts this chain at its first link.
After spreads between the sovereigns of the union have widened, the central bank
announces that it will not let the spread of the $D$ sovereign over the $F$
sovereign exceed a ceiling, and that it will buy $D$ bonds in whatever quantity
it takes to enforce it. The design follows the ECB's Outright Monetary
Transactions (2012) and Transmission Protection Instrument (2022): both are
conditional on the state of the bond market and unlimited in size.

We model the announcement as a permanent change of policy regime that every
agent understands. The comparison is between two economies that are identical
except for the backstop: in one, the purchase rule below is part of the
equilibrium and enters every expectation; in the other, it does not exist. Two
effects follow, and the model separates them. The \emph{announcement effect} is
the difference between the two economies at their stochastic rest points, where
the backstop is in place but nothing is bought. The \emph{intervention effect} is
the difference between their responses to the same rise in $\pi^{d}_t$, where it
may be used.

\paragraph{The objective and the purchase rule:}
The model is real, so we abstract from the central bank's price-stability
mandate and from inflation dynamics, and state its objective in the bond market
alone, in terms of the sovereign spread $sp_t=y(q^{D}_t)-y(q^{F}_t)$ of
\eqref{eq:spread-def}, the gap between the two flow yields \eqref{eq:hm-yield}.
The central bank announces a cap $\bar c$ on that spread and enforces it with the
smallest intervention that does so:
\begin{equation}
\min_{m^{cb}_t\,\ge\,0}\;m^{cb}_t
\qquad\text{s.t.}\qquad
y\bigl(q^{D}_t\bigr)-y\bigl(q^{F}_t\bigr)\;\le\;\bar c
\quad\text{if } d_t=0,
\qquad\qquad
m^{cb}_t=0\quad\text{if } d_t=1,
\label{eq:cb-objective}
\end{equation}
where $m^{cb}_t$ is the quantity of $D$ bonds bought at $t$, in the issuer's
per-capita units. Because \eqref{eq:hm-yield} is decreasing in $q$, the cap is a
floor on the price of the $D$ bond:
\begin{equation}
q^{D}_t\;\ge\;\underline q_t\;\equiv\;\frac{\delta_b}{\delta_b+y\bigl(q^{F}_t\bigr)+\bar c}.
\label{eq:tpi-floor}
\end{equation}
Purchases raise the market price of the bond (a property we verify below), so
the solution of \eqref{eq:cb-objective} is the complementarity condition
\begin{equation}
m^{cb}_t\;\ge\;0,
\qquad
q^{D}_t-\underline q_t\;\ge\;0,
\qquad
m^{cb}_t\,\bigl(q^{D}_t-\underline q_t\bigr)\;=\;0
\qquad\text{in } d_t=0,
\label{eq:tpi-kkt}
\end{equation}
which is also the decision rule under which the central bank operates. As long as
the market clears the $D$ bond above the floor, the central bank does nothing;
when the market would clear below it, the central bank buys exactly the quantity
that holds the price at $\underline q_t$, and no more. The problem is
deliberately a constrained one: the central bank does not minimise the spread.
An objective that rewarded a lower spread without limit would have the central
bank buying in every state, for as long as one more purchase lowered the spread,
and would turn the backstop into a permanent holding of the periphery's debt.
The cap is what makes the policy a backstop: inactive inside the band,
unlimited at its edge.

Four features of \eqref{eq:cb-objective}--\eqref{eq:tpi-kkt} matter for what
follows. First, the backstop operates only in the no-default regime: it insures
bondholders against the \emph{price} of default risk, not against default
itself. Second, purchases are held to maturity: $m^{cb}_t\ge0$, and bonds are
never sold back. Third, the floor moves with the safe yield. When a flight to
safety lowers $y(q^{F}_t)$, $\underline q_t$ rises and the central bank must buy
more, precisely in the states where buying is dearest. Fourth, the cap is set in
the flow-yield metric \eqref{eq:hm-yield}, the same metric in which we report
spreads; in the baseline $\bar c$ is 200 basis points a year.

\paragraph{How purchases are paid for:}
The central bank pays for its bonds with a riskless one-period claim on itself,
held by the $D$ banks and remunerated at the $D$ deposit rate: interest-bearing
reserves, in effect. Let $M^{cb}_t$ be the stock of $D$ bonds the central bank
carries into period $t$. Its book, the reserves outstanding at the end of $t$ and
the obligation they create evolve as
\begin{equation}
M^{cb}_{t+1}=(1-\delta_b)\,h_t\,M^{cb}_t+m^{cb}_t,
\qquad
\mathrm{res}_t=q^{D}_t\,M^{cb}_{t+1},
\qquad
O^{cb}_{t+1}=R_{D,t}\,\mathrm{res}_t .
\label{eq:cb-book}
\end{equation}
The central bank holds the bond on the same terms as any other investor: its
book decays at the maturity rate and is written down by the same survival factor
$h_t$ in default (it ranks \emph{pari passu}). Its profit in period $t$ is the
payoff on the book less the obligation on the reserves issued at $t-1$,
\begin{equation}
\Pi^{cb}_t=\Xi^{D}_t\,M^{cb}_t-O^{cb}_t,
\qquad
\Xi^{D}_t=h_t\bigl[\delta_b+(1-\delta_b)\,q^{D}_t\bigr],
\label{eq:cb-pl}
\end{equation}
and it is remitted in full every period, so the central bank carries no net
worth. Given \eqref{eq:cb-book}, its cash flows
\begin{equation*}
\underbrace{\delta_b\,h_t\,M^{cb}_t}_{\text{coupons}}
+\underbrace{\mathrm{res}_t}_{\text{new reserves}}
\;=\;
\underbrace{O^{cb}_t}_{\text{reserves repaid}}
+\underbrace{\Pi^{cb}_t}_{\text{remittance}}
+\underbrace{q^{D}_t\,m^{cb}_t}_{\text{purchases}}
\end{equation*}
balance identically. The profit is shared by the euro-area capital key: a share
$\kappa_D$ goes to the $D$ treasury and $1-\kappa_D$ to the $F$ treasury.

\paragraph{The two treasuries:}
The $D$ budget constraint \eqref{eq:govt-budget} gains the remittance,
\begin{equation}
G_D+\delta_b\,h_t\,B_{D,t}
\;=\;T^{\tau}_{D,t}+\kappa_D\,\Pi^{cb}_t
+q^{D}_t\Bigl[B_{D,t+1}-(1-\delta_b)\,h_t\,B_{D,t}\Bigr],
\qquad
B_{D,t}=b^{D}_{D,t}+b^{D}_{F,t}+M^{cb}_t,
\label{eq:govt-budget-cb}
\end{equation}
where $B_{D,t}$ is now the gross stock, including the central bank's holdings.
Coupons and the Bohn rule \eqref{eq:bohn} apply to the gross stock: the
government pays the central bank as it pays any other holder. The $F$ stock is
fixed at $\bar B_\F$, so $F$'s share of the profit lowers $F$ taxes, expressed in
$F$-goods per $F$ capita:
\begin{equation}
T^{\tau}_{F,t}\;=\;\bar T^{\tau}_{F}-\frac{(1-\kappa_D)\,\Pi^{cb}_t}{\varkappa\,p_t}.
\label{eq:tax-F-cb}
\end{equation}
The remittance is how the backstop redistributes across the union. Along a path
without default, the central bank earns the excess return on the bonds it bought,
and $\Pi^{cb}_t>0$. In default it takes the haircut on its book, $\Pi^{cb}_t<0$,
and the loss is split by the capital key. With $\kappa_D=0.071$, 92.9\% of either
goes to $F$ households.

\paragraph{The $D$ bank's balance sheet:}
The reserves are an asset of the $D$ bank and the obligation on them is part of
its payoff, so the balance sheet \eqref{eq:balance-sheet} and the payoff
\eqref{eq:asset-payoff} become
\begin{align}
\mathcal{A}_{D,t}
&=Q_{D,t}K_{D,t+1}+q^{D}_t\,b^{D}_{D,t+1}+p_t\,q^{F}_t\,b^{F}_{D,t+1}+L_{D,t}
 +\mathrm{res}_t
 \;=\;n_{D,t}+\mathrm{dep}_{D,t},
\label{eq:balance-sheet-cb}\\[3pt]
\mathcal{X}_{D,t+1}
&=\bigl[mpk_{D,t+1}+(1-\delta_D)Q_{D,t+1}\bigr]K_{D,t+1}
 +\Xi^{D}_{t+1}\,b^{D}_{D,t+1}+p_{t+1}\,\Xi^{F}_{t+1}\,b^{F}_{D,t+1}
 +O^{cb}_{t+1},
\label{eq:asset-payoff-cb}
\end{align}
and the market for the $D$ bond clears net of the central bank's book:
\begin{equation}
b^{D}_{D,t+1}=B_{D,t+1}-b^{D}_{F,t+1}-M^{cb}_{t+1}.
\label{eq:bond-clearing-cb}
\end{equation}
Net worth \eqref{eq:nw-lom}, deposits, the carried obligation \eqref{eq:P-state},
the entrant transfer and bank dividends \eqref{eq:bank-div} keep their formulas,
applied to the enlarged $\mathcal{A}_{D,t}$ and $\mathcal{X}_{D,t+1}$.

This involves two modelling choices. First, reserves are not a portfolio choice:
the purchase creates them and the $D$ banks must hold all of them, so there is no
first-order condition for $\mathrm{res}_t$. Second, reserves enter the incentive
constraint \eqref{eq:IC} at the same divertibility $\lambda_D$ as every other
asset. This is the conservative choice. It denies the backstop any mechanical
relief of the constraint, so whatever the policy achieves must come through
prices and future payoffs. This is also where it differs from a collateralised
lending facility, whose credit cannot be diverted and counts as equity in the
constraint.

\paragraph{The bank may sell its whole book:}
Without the backstop, the bank's bond condition \eqref{eq:bond-pricing} always
holds with equality: the price falls until the bank is willing to hold the
residual supply. With the backstop, the price can be held above the bank's own
valuation, and the bank then sells to the central bank. Because its holdings are
bounded below by zero, its first-order condition becomes a complementarity pair:
\begin{equation}
b^{D}_{D,t+1}\ge0,
\qquad
\mathcal{F}_{D,t}\equiv
\E_t\bigl[\Omega_{D,t,t+1}\,\Xi^{D}_{t+1}\bigr]
-q^{D}_t\Bigl(\E_t\bigl[\Omega_{D,t,t+1}\bigr]R_{D,t}+\lambda_D\,\mu_{D,t}\Bigr)\le0,
\qquad
b^{D}_{D,t+1}\,\mathcal{F}_{D,t}=0 .
\label{eq:bank-bond-kkt}
\end{equation}
Together, \eqref{eq:tpi-kkt} and \eqref{eq:bank-bond-kkt} produce three cases in
the no-default regime:
\begin{enumerate}
\item \emph{Floor slack} ($q^{D}_t>\underline q_t$). Nothing is bought,
  $m^{cb}_t=0$, and $q^{D}_t$ is given by \eqref{eq:bond-pricing}, with
  expectations formed in the backstop economy.
\item \emph{Floor binding, interior} ($q^{D}_t=\underline q_t$,
  $b^{D}_{D,t+1}>0$). The bank's condition holds at the floor price,
  $\mathcal{F}_{D,t}=0$, and $m^{cb}_t>0$ is the purchase at which the bank values
  the bond at exactly $\underline q_t$.
\item \emph{Floor binding, corner} ($q^{D}_t=\underline q_t$,
  $b^{D}_{D,t+1}=0$). The bank values the bond below the floor even when it holds
  none, $\mathcal{F}_{D,t}<0$, and the central bank holds every bond the $F$ bank
  does not: $M^{cb}_{t+1}=B_{D,t+1}-b^{D}_{F,t+1}$.
\end{enumerate}
In the default regime nothing is bought, and the price is set by
\eqref{eq:bank-bond-kkt} alone. The $F$ bank's condition \eqref{eq:bondD-foreign}
is unchanged: its adjustment cost makes its holdings fall as the price rises, so
at the floor it sells to the central bank as well.

Case 2 deserves a remark. The banker's value function is linear in net worth, so
\eqref{eq:bank-bond-kkt} contains no quantity. A purchase can move the bank's
valuation only through equilibrium objects: through $\mu_{D,t}$, which the swap
below leaves unchanged on impact, and through the kernel $\Omega_{D,t,t+1}$,
which depends on the next-period state
$(b^{D}_{D,t+1},M^{cb}_{t+1},O^{cb}_{t+1})$. The inverse demand for the $D$
bond,
$q^{D}_t=q^{D}\bigl(B_{D,t+1}-M^{cb}_{t+1};\,n_{D,t},\mu_{D,t},\Omega_{D},\pi^{d}_t\bigr)$,
is therefore a general-equilibrium object, and its slope in $M^{cb}$ is a result
rather than an assumption. We deliberately add no home-bias or preferred-habitat
cost to the bank's portfolio problem. Such a cost would put a quantity into
\eqref{eq:bank-bond-kkt} by assumption, and the price effect of purchases would
then be that assumption.

\paragraph{A purchase is a swap today and a risk transfer tomorrow:}
Consider a marginal purchase $dm$ at price $q^{D}_t$, holding prices and the
bank's other choices fixed. The bank gives up $dm$ bonds,
$db^{D}_{D,t+1}=-dm$, and receives reserves $d\,\mathrm{res}_t=q^{D}_t\,dm$. Its
assets, the divertible base $\lambda_D\mathcal{A}_{D,t}$, deposits and the carried
obligation $P_{D,t+1}$ are all unchanged, and so, by \eqref{eq:mu-closed}, is
$\mu_{D,t}$. Within the period, the purchase does nothing to the bank. Next
period, in regime $d'$, the bank is owed $R_{D,t}\,q^{D}_t\,dm$ on its reserves in
place of the bonds' payoff $\Xi^{D}_{t+1}(d')\,dm$. By \eqref{eq:asset-payoff-cb},
\eqref{eq:cb-pl} and \eqref{eq:nw-lom},
\begin{equation}
dn_{D,t+1}(d')=\bigl[(1-f_D)+\omega^{f}_{D}\,\kappa_D\bigr]
\bigl(R_{D,t}\,q^{D}_t-\Xi^{D}_{t+1}(d')\bigr)\,dm,
\qquad
d\Pi^{cb}_{t+1}(d')=-\bigl(R_{D,t}\,q^{D}_t-\Xi^{D}_{t+1}(d')\bigr)\,dm .
\label{eq:risk-transfer}
\end{equation}
The term $\omega^{f}_{D}\kappa_D$ arises because $D$'s share of the central bank's
profit changes $D$'s new issuance by $-\kappa_D\,d\Pi^{cb}_{t+1}/q^{D}_{t+1}$; the
$D$ bank absorbs the new bonds, and they enter the entrants' transfer through
$\mathcal{A}_{D,t+1}$. Without default,
$\Xi^{D}_{t+1}(0)=\delta_b+(1-\delta_b)q^{D}_{t+1}$ exceeds $R_{D,t}q^{D}_t$ by
the bond's excess return, and the bank gives up a little of it. In default,
$\Xi^{D}_{t+1}(1)=\varrho_D[\delta_b+(1-\delta_b)q^{D}_{t+1}]$ falls far short
of $R_{D,t}q^{D}_t$, and the bank is insured: it holds a claim that pays in full
where the bond it sold is written down by $1-\varrho_D$. In the calibrated model
the net-worth gain in the default branch is two orders of magnitude larger than
the cost in the no-default branch.

The insurance reaches today's price through the banker's kernel. Higher net
worth in the default branch lowers $\mu^{(1)}_{D,t+1}$, and with it the franchise
value $\varphi^{(1)}_{D,t+1}$ in \eqref{eq:phi-recursion}, which lowers
$\Omega^{(1)}_{D,t,t+1}$ in \eqref{eq:branch-sdf}. The covariance between the
kernel and the bond's payoff in \eqref{eq:bond-decomposition} shrinks, the risk
premium falls and $q^{D}_t$ rises. A franchise channel works in the opposite
direction: a safer future lowers $\varphi_{D,t+1}$, hence
$\E_t[\Omega_{D,t,t+1}]$, which by \eqref{eq:mu-closed} tightens today's
constraint. Which channel dominates is a quantitative question. At the calibrated
rest point $\E_t[\Omega_{D,t,t+1}]$ is higher in the backstop economy, so the
franchise channel does not dominate.

\paragraph{Credibility and the probability of intervention:}
Financial intermediaries price the backstop through their assessment of how
likely it is to be used. In the model the rule is fully credible: agents know the
purchase function $m^{cb}(\mathcal{S})$ and know that
$q^{D}(0,\mathcal{S})\ge\underline q(\mathcal{S})$ at every state. There is
therefore no exogenous probability of intervention to assess. The probability
that the central bank intervenes next period is nonetheless well defined, and it
is endogenous. Under \eqref{eq:expectation},
\begin{equation}
\Pr_t\bigl(m^{cb}_{t+1}>0\bigr)
=\bigl(1-\pi^{d}_t\bigr)\,
\E^{s}_t\Bigl[\mathbb{1}\bigl\{\mathcal{S}_{t+1}\in\mathcal{S}^{\rm bind}\bigr\}\Bigr],
\label{eq:prob-intervention}
\end{equation}
where $\mathcal{S}^{\rm bind}$ is the set of states at which the floor binds. A
higher $s_t$ raises this probability through the persistence of
\eqref{eq:s-process}, which moves next period's state towards
$\mathcal{S}^{\rm bind}$. At the same time it lowers the weight $1-\pi^{d}_t$ on
the regime in which the central bank buys at all.

It is useful to write the next-period no-default price as
$q^{D}_{t+1}(0)=\max\{\tilde q_{t+1},\underline q_{t+1}\}$, where
$\tilde q_{t+1}$ is the price at which the market would clear at $t+1$ with no
purchase at that date but with the backstop in place thereafter. The expected
no-default payoff then splits into
\begin{equation}
\E^{s}_t\bigl[\Xi^{D}_{t+1}(0)\bigr]
=\delta_b+(1-\delta_b)\,\E^{s}_t\bigl[\tilde q_{t+1}\bigr]
+\underbrace{(1-\delta_b)\,\E^{s}_t\Bigl[\bigl(\underline q_{t+1}-\tilde q_{t+1}\bigr)^{+}\Bigr]}_{\text{the backstop's put}} .
\label{eq:put}
\end{equation}
The last term is the value of a put option on the $D$ bond, written by the
central bank, struck at the floor and exercisable only if the sovereign has not
defaulted. It enters today's price through the numerator of
\eqref{eq:bond-pricing}, weighted by the kernel in the no-default branch. It does
so at every state, including states where the floor is far from binding today and
the put is out of the money. This is how the announcement alone moves the market.

Credibility also lowers the cost of the backstop. The more the market prices the
put, the higher $\tilde q$ is relative to the price in an economy without the
backstop, and the less the central bank must buy when the floor binds. The
purchase only has to close the gap between the floor and the backstop economy's
own market price, not the gap to the price that would prevail without the
backstop. The accounting cost is deferred as well. Along any path without default
the central bank earns the bond's excess return; it makes losses only if the
sovereign defaults while the book is still held.

\paragraph{Time consistency:}
The objective \eqref{eq:cb-objective} is static and the rule \eqref{eq:tpi-kkt}
is Markov, a function of the current state alone. The purchase the rule
prescribes at any date is therefore the one a central bank re-solving
\eqref{eq:cb-objective} at that date would choose. Following the rule requires no
commitment beyond the objective itself, and in this sense the backstop is time
consistent. This is a property of the objective we assume, not a result. A
central bank that also weighed the expected loss on its book,
$-\pi^{d}_t\,\E^{s}_t\bigl[\Pi^{cb}_{t+1}(1)\bigr]$, would be tempted to renege in
exactly the high-risk states where the floor binds, and the announcement would
then be priced with a probability of reneging. We abstract from that temptation.
If the announcement alone kept the spread below the cap, the commitment would
never be tested. In the model this holds in normal times, since nothing is bought
at the rest point, but not in a crisis, when the floor binds and the central bank
buys.

\paragraph{State and equilibrium:}
The book and the obligation on reserves are two new states, both zero at the
steady state. $O^{cb}_t$ must be carried separately from $M^{cb}_t$ because it
embeds the price at which the reserves were issued,
$O^{cb}_t=R_{D,t-1}\,q^{D}_{t-1}\,M^{cb}_t$, which nothing else in the state
records. The backstop is a rule, not a random event, so expectations are taken
under \eqref{eq:expectation} over the default regime alone. At each state and
regime the equilibrium has one more unknown, the purchase $m^{cb}_t$, and one
more condition, \eqref{eq:tpi-kkt}; the $D$ bank's $D$-bond condition takes the
complementarity form \eqref{eq:bank-bond-kkt}. The deterministic steady state is
unchanged ($M^{cb}=O^{cb}=\mathrm{res}=\Pi^{cb}=0$), and with the backstop
switched off the model is exactly the one of the preceding sections.

\paragraph{Computation:}
The complementarity conditions \eqref{eq:tpi-kkt} and \eqref{eq:bank-bond-kkt}
put kinks into the decision rules. We solve the purchase pair in the form of
\cite{garcia1981pathways}, in which a single smooth variable $\chi_t$ carries both
sides,
\begin{equation}
m^{cb}_t=\bar B_\D\,\max\{\chi_t,0\},
\qquad
q^{D}_t=\underline q_t+\bar q^{D}\,\max\{-\chi_t,0\},
\label{eq:gz}
\end{equation}
so that \eqref{eq:tpi-kkt} holds by construction and the stored rule for
$\chi_t$ is smooth across the floor. The bank's corner \eqref{eq:bank-bond-kkt} is
imposed as a \cite{fischer1992special} pair. Both are smoothed at
$\varepsilon=10^{-4}$. The cap is walked down in 50 bp steps from 700 bp, where
it binds nowhere on the state space, to its target. Details are in the
computational appendix.


% -----------------------------------------------------------------------------
%  RESULTS. Kept here so the section reads end to end; move to the results
%  section if preferred. Source and caveat: see the header of this file.
% -----------------------------------------------------------------------------

\paragraph{The backstop at work:}
At the baseline cap of 200 bp the floor binds at 43 of the 115 no-default
collocation nodes. These are every node at the highest risk level on the grid
($\pi^{d}=4.8\%$ a quarter) and most of those at the next. At 21 of them, all at
the highest risk level, the bank sells its entire book to the central bank (case
3).

\begin{table}[ht]
\centering
\caption{The announcement effect: stochastic rest points, cap 200 bp}
\label{tab:tpi-rest}
\begin{tabular}{lccc}
\hline
 & No backstop & Backstop & Difference\\
\hline
Sovereign spread, bp/yr          & 55.8   & 51.6   & $-4.3$\\
Credit spread, bp/yr             & 88.7   & 52.6   & $-36.1$\\
$\mu_D$ (IC multiplier)          & 0.0110 & 0.0066 & $-0.0045$\\
$\E[\Omega_D]$                   & 1.160  & 1.166  & $+0.53\%$\\
$n_D$ (bank net worth)           & 1.939  & 1.948  & $+0.44\%$\\
$Y_D$ (output)                   & 0.9994 & 1.0012 & $+0.18\%$\\
$I_D$ (investment)               & 0.2046 & 0.2057 & $+0.56\%$\\
Purchases, \% of $\bar B_\D$     & 0      & 0      & 0\\
\hline
\end{tabular}
\end{table}

The announcement works mainly through the banks' balance sheets, not through
the sovereign spread (\cref{tab:tpi-rest}). At the rest point the floor is far
out of the money: the spread is 56 bp against a cap of 200 bp. The bond price
rises by only 0.24\% and the sovereign spread falls by 4 bp. Yet the multiplier
$\mu_D$ falls by two-fifths and the credit spread by 36 bp, because both terms in
the numerator of \eqref{eq:mu-closed} rise: net worth by 0.44\% and
$\E[\Omega_D]$ by 0.53\%. Output is 0.18\% higher and investment 0.56\% higher,
and nothing is bought.

\begin{table}[ht]
\centering
\caption{The intervention effect: $\pi^{d}_t$ rises to 2\% a quarter, cap 200 bp}
\label{tab:tpi-shock}
\begin{tabular}{lcc}
\hline
 & No backstop & Backstop\\
\hline
Output, impact                              & $-0.128\%$ & $-0.076\%$\\
Output, trough                              & $-0.128\%$ (q0) & $-0.149\%$ (q6)\\
Consumption, impact                         & $-0.02\%$ & $-0.24\%$\\
Sovereign spread, impact, bp/yr             & 303 & 200\\
Credit spread, impact, bp/yr                & 180 & 31\\
Bank net worth, impact                      & $-4.19\%$ & $-1.45\%$\\
$q^{D}$, impact                             & $-8.97\%$ & $-5.68\%$\\
Purchases, impact, \% of $\bar B_\D$        & --- & 44.3\\
\hline
\multicolumn{3}{p{0.75\linewidth}}{\footnotesize Quantities are deviations
from each economy's own unshocked path, starting at its own rest point; spreads
are levels. Both paths clear the period map exactly every quarter.}
\end{tabular}
\end{table}

When $\pi^{d}_t$ jumps to 2\% a quarter, the backstop fires once
(\cref{tab:tpi-shock}). The central bank buys 44\% of the steady-state stock on
impact and holds the spread at the cap, 200 bp against 303 bp without the
backstop. From the next quarter on the market spread is below the cap, nothing
more is bought, and the book runs off at the maturity rate. The purchase is
interior (case 2): the $D$ bank keeps part of its book. The impact fall in output
is 40\% smaller, bank net worth falls by 1.4\% instead of 4.2\%, and the credit
spread is 31 bp on impact instead of 180 bp.

The intervention works mainly through the continuation value of the bond.
Because the bank's condition holds with equality at the headline shock, the bond
price factors exactly into
\begin{equation}
\log q^{D}_t
=-\log R_{D,t}
+\log\E^{s}_t\bigl[\Xi^{D}_{t+1}(0)\bigr]
+\log\frac{\E_t\bigl[\Xi^{D}_{t+1}\bigr]}{\E^{s}_t\bigl[\Xi^{D}_{t+1}(0)\bigr]}
+\log\frac{\E_t\bigl[\Omega_{D}\Xi^{D}_{t+1}\bigr]}
          {\E_t[\Omega_{D}]\,\E_t\bigl[\Xi^{D}_{t+1}\bigr]}
+\log\frac{\E_t[\Omega_{D}]\,R_{D,t}}{\E_t[\Omega_{D}]\,R_{D,t}+\lambda_D\mu_{D,t}},
\label{eq:bond-log-decomp}
\end{equation}
where $\Xi^{D}_{t+1}(0)$ is the payoff in the no-default regime: discounting, the
continuation price, the expected default loss, the risk premium and the liquidity
premium, the exact counterpart of \eqref{eq:bond-decomposition}. The backstop
economy prices the bond 3.89\% higher on impact. Of that, the continuation price
contributes 3.31 percentage points, the liquidity premium 0.38, the risk premium
0.31 and the expected default loss 0.05, against $-0.16$ from discounting. The
purchase itself does little. Holding the decision rules fixed, a one-off purchase
of 30\% of the stock raises $q^{D}_t$ by 0.61\% at the headline shock and by
0.24\% at the rest point. The backstop supports the price through what it
promises for tomorrow, the floor under the bond's resale value in \eqref{eq:put},
rather than through what it buys today.

The downturn is shallower on impact but longer. Output troughs at $-0.149\%$ in
quarter 6, whereas without the backstop it turns positive within a year. An
impulse response conditions on no default, and along such a path the risk
transfer \eqref{eq:risk-transfer} runs in its no-default direction. The $D$ banks
sold 44\% of the stock at the floor and hold reserves instead. The bonds' excess
return, which in the economy without the backstop rebuilds bank net worth within
a year, therefore accrues to the central bank, and 92.9\% of it goes to $F$
households through \eqref{eq:tax-F-cb}. $D$ bank net worth stays depressed, the
credit spread climbs to 83 bp by quarter 6, and $D$ consumption falls by 0.24\% on
impact and 0.34\% at its trough. The impulse response therefore shows what the
insurance costs and none of what it pays. Its benefit lies in the default branch,
where the central bank, and through the capital key mostly $F$, takes the haircut
in place of the $D$ banks. A welfare comparison needs that branch.


% =============================================================================
%  INSERT: replaces eq. (state) and the RCE definition in
%          \subsection{Recursive Competitive Equilibrium} of 02-model.tex.
%          The surrounding prose of that subsection is unchanged, except that
%          "three of the four carried sovereign positions" now also needs the
%          book: B_{D,t} = b^D_{D,t} + b^D_{F,t} + M^{cb}_t.
% =============================================================================

\begin{equation}
\mathcal{S}_t=\Bigl[
K_{\D,t},\;K_{\F,t},\;
P_{\D,t},\;P_{\F,t},\;
b^{D}_{D,t},\;b^{D}_{F,t},\;b^{F}_{D,t},\;
V_t,\;s_t,\;Z_{\D,t},\;
M^{cb}_t,\;O^{cb}_t\Bigr],
\label{eq:state}
\end{equation}

\begin{definition}[Recursive competitive equilibrium]
\label{def:rce}
A recursive competitive equilibrium is a set of decision rules, one for each
regime $d\in\{0,1\}$,
\[
\bigl\{N_X,\;K'_X,\;r_X,\;A_X,\;C_X,\;\varphi_X,\;r^{wc}_X,\;
p,\;q^{D},\;q^{F},\;b^{D}_{F}{}',\;b^{F}_{D}{}',\;m^{cb}\bigr\}(d,\mathcal{S}),
\]
the implied multipliers $\mu_X$ from \eqref{eq:mu-closed}, and a law of motion
for $\mathcal{S}$ generated by \eqref{eq:capital-lom}, \eqref{eq:P-state},
\eqref{eq:govt-budget-cb}, \eqref{eq:V-lom}, \eqref{eq:s-process},
\eqref{eq:z-process} and \eqref{eq:cb-book}, such that at every
$(d,\mathcal{S})$: households satisfy \eqref{eq:agg-budget}--\eqref{eq:agg-euler};
firms and capital producers satisfy \eqref{eq:labour-demand}--\eqref{eq:tobins-q};
intermediaries satisfy \eqref{eq:foc-general} for every asset class, with the
$D$ bank's $D$-bond condition in the form \eqref{eq:bank-bond-kkt}, together with
\eqref{eq:mu-closed}; the governments satisfy \eqref{eq:govt-budget-cb},
\eqref{eq:tax-F-cb} and \eqref{eq:bohn}; the central bank follows
\eqref{eq:tpi-kkt}; expectations are taken under \eqref{eq:expectation}; and the
goods, sovereign \eqref{eq:bond-clearing-cb}, deposit, capital and labour markets
clear.
\end{definition}
